% Paper 6 (GRPO-Registry) worked examples
\section{Worked Examples}
\label{sec:p6-worked}
Both examples below are verbatim CLI runs on the shipped entries.

\subsection{Example 1: \texttt{stackdiff} on the ``DAPO'' Pair}
\label{sec:p6-worked-dapo}
The registry contains two stacks claiming the label \texttt{dapo}: the open
Colab trainer arm and the managed-runtime arm. Diffing them:

\begin{lstlisting}[basicstyle=\ttfamily\scriptsize, frame=single,
  caption={\texttt{stackdiff colab-open_dapo_e3
  tinker_dapo_qwen3.5-4b_gsm8k} (abridged: R2 backend lines omitted).},
  label={lst:p6-stackdiff}]
stackdiff colab-open_dapo_e3 vs tinker_dapo_qwen3.5-4b_gsm8k
  shared claimed label: 'dapo'
  [R5] loss_form.importance_ratio_level: 'token' -> None
  [R5] reference_kl.kl_beta: 0.0 -> None
  [R3] reference_kl.reference_policy: False -> True
  [R4] delta delta_dapo:clip_higher: implemented -> surrogate
  [R4] delta delta_dapo:dynamic_sampling: implemented -> absent
  [R5] delta delta_dapo:kl_removed: implemented -> unknown
  [R5] delta delta_dapo:token_level_loss: implemented -> unknown
verdict: R5 unauditable: a differing item is unreported (null)
         on at least one side
verdict: comparisons across this pair are at risk of a LABEL FLIP;
         do not attribute outcome differences to the shared label.
\end{lstlisting}

The verdict is R5 with R4 findings inside it: DAPO's dynamic sampling---the
component that mechanically drives batch ZVF to zero by filtering
zero-variance groups---is \texttt{implemented} on one side and \texttt{absent}
on the other, while the KL and loss-aggregation components cannot even be
compared because the managed side leaves them unreported. This is a
\emph{pre-run} verdict computed from manifests alone. The program's telemetry
subsequently confirmed it empirically: mean ZVF $0.00$ on the open arm versus
$0.58$ on the managed arm under the identical label (toy scale on the open
side, hence directional---but the direction is the one the manifest diff
predicted, and the mechanism is definitional). The general lesson matches the
backend-swap audit, where a same-label, R2-level swap moved final reward
$85.6\% \to 5.0\%$: rankings inherit the stack, and \texttt{stackdiff} makes
the inheritance legible before compute is spent.

\subsection{Example 2: ``Which Stacks Report KL Handling?''}
\label{sec:p6-worked-kl}
Item 2 (reference policy + KL) is the canonical scattered-in-appendices
field: defaults differ silently across frameworks ($\beta{=}0.04$ in TRL and
veRL versus $0.02$ in OpenRLHF, per their own dumps), and DAPO removes the
term outright. The registry answers in one call:

\begin{lstlisting}[basicstyle=\ttfamily\scriptsize, frame=single,
  caption={\texttt{query --item reference_kl} on the seed entries
  (grouped).}, label={lst:p6-klquery}]
colab-open_{grpo,drgrpo,dapo,grpo-adaptiveg}   reference_kl=FULL (1.00)
trl_grpo_qwen3-8b_gsm8k                        reference_kl=FULL (1.00)
verl_grpo_qwen3-8b_gsm8k                       reference_kl=FULL (1.00)
openrlhf_grpo_qwen3-8b_gsm8k                   reference_kl=FULL (1.00)
tinker_{grpo(8B),grpo,dapo,drgrpo,gspo}        reference_kl=PARTIAL (0.33)
\end{lstlisting}

Seven of twelve stacks fully report KL handling. The instructive cases are the
two ends: the Colab arms score \texttt{FULL} by \emph{explicitly declaring the
absence} of a reference policy ($\beta{=}0.0$, estimator \texttt{none})---a
reported negative is a first-class answer---while all five managed-runtime
entries are \texttt{PARTIAL} because \texttt{kl.beta} and the estimator are
\texttt{managed_by_tinker} nulls. A meta-analyst asking ``does KL
regularization explain the variant ranking?'' learns, in one query, exactly
which stacks can enter the regression and which must be excluded as
unauditable on that axis.
